% Generated from lessons/0009-4-arbitrary-series.html; edit the HTML source, then rerun the converter.
\section{任意项级数}
\lessonmark{任意项级数}

正项级数只需要研究“部分和有没有上界”；任意项级数的难点在于正负项会互相抵消。本节的主线是：先用 Cauchy 收敛原理抓住本质，再学习交错级数、Abel 与 Dirichlet 判别法，最后区分绝对收敛、条件收敛、重排和乘法。

\subsection*{本节主线}

\begin{conceptbox}
\textbf{先看尾和}
任意项级数不再有单调部分和，必须用 Cauchy 收敛原理控制任意一段尾和。
\end{conceptbox}

\begin{conceptbox}
\textbf{再看抵消}
交错级数、三角级数等靠正负抵消收敛，重点不是每项绝对值够小，而是部分和能被控制。
\end{conceptbox}

\begin{conceptbox}
\textbf{最后分类型}
绝对收敛最稳，条件收敛最敏感：重排可能改和值，乘法也需要额外条件。
\end{conceptbox}

\begin{notebox}
考试顺序：先看 \(\sum |x_n|\)；若绝对收敛，原级数必收敛。若不绝对收敛，再看是否是 Leibniz、Dirichlet 或 Abel 型。
\end{notebox}

\subsection*{一、任意项级数的 Cauchy 收敛原理}

一个级数如果只有有限个正项或有限个负项，就可以转化为正项级数讨论。真正新的情况是：既有无穷多个正项，又有无穷多个负项。

\subsubsection*{定理 9.4.1：级数的 Cauchy 收敛原理}

级数 \(\sum_{n=1}^{\infty}x_n\) 收敛的充要条件是：对任意 \(\varepsilon>0\)，存在正整数 \(N\)，使得当 \(m>n>N\) 时，

\[
      \left|x_{n+1}+x_{n+2}+\cdots+x_m\right|
      =
      \left|\sum_{k=n+1}^{m}x_k\right|
      <\varepsilon.
    \]

等价地，对任意 \(\varepsilon>0\)，存在 \(N\)，使得对一切 \(n>N\) 和一切正整数 \(p\)，

\[
      \left|x_{n+1}+x_{n+2}+\cdots+x_{n+p}\right|
      =
      \left|\sum_{k=1}^{p}x_{n+k}\right|
      <\varepsilon.
    \]

取 \(p=1\)，立刻得到必要条件：

\[
      \sum_{n=1}^{\infty}x_n\text{ 收敛}\quad\Longrightarrow\quad
      \lim_{n\to\infty}x_n=0.
    \]

\begin{notebox}
注意：一般项趋于 \(0\) 仍然只是必要条件，不是充分条件。本节很多题就是在考“趋于 0 之外，还靠什么收敛”。
\end{notebox}

\subsection*{二、Leibniz 级数与交错判别法}

\subsubsection*{定义 9.4.1：交错级数与 Leibniz 级数}

若 \(u_n>0\)，形如

\[
      \sum_{n=1}^{\infty}(-1)^{n+1}u_n
    \]

的级数称为交错级数。若进一步 \(\{u_n\}\) 单调减少且收敛于 \(0\)，则称它为 Leibniz 级数。

\subsubsection*{定理 9.4.2：Leibniz 判别法}

Leibniz 级数必定收敛。也就是说，若

\[
      u_1\ge u_2\ge \cdots \ge 0,\qquad \lim_{n\to\infty}u_n=0,
    \]

则

\[
      \sum_{n=1}^{\infty}(-1)^{n+1}u_n
    \]

收敛。

\subsubsection*{证明思路}

关键不是把正项和负项分开，而是估计任意一段尾和。因为 \(u_n\) 单调减少，尾和

\[
      u_{n+1}-u_{n+2}+u_{n+3}-\cdots+(-1)^{p+1}u_{n+p}
    \]

的绝对值总不超过第一项 \(u_{n+1}\)。既然 \(u_{n+1}\to0\)，尾和就能任意小，于是由 Cauchy 收敛原理得收敛。

\subsubsection*{余项估计}

若

\[
      r_n=\sum_{k=n+1}^{\infty}(-1)^{k+1}u_k,
    \]

则

\[
      |r_n|\le u_{n+1}.
    \]

这在考试里很有用：不仅能判收敛，还能估计用前 \(n\) 项近似时的误差。

\begin{conceptbox}
\textbf{典型例子}
\(\sum_{n=1}^{\infty}\frac{(-1)^{n+1}}{n}\) 收敛，但不绝对收敛。
\end{conceptbox}

\begin{conceptbox}
\textbf{容易漏掉的条件}
只知道 \(u_n\to0\) 不够，还需要 \(u_n\) 最终单调减少，或换用 Dirichlet 判别法。
\end{conceptbox}

\subsection*{三、Abel 变换、Dirichlet 判别法与 Abel 判别法}

\subsubsection*{引理 9.4.1：Abel 变换}

设 \(\{a_n\}\)、\(\{b_n\}\) 是两数列，记

\[
      B_k=\sum_{i=1}^{k}b_i.
    \]

则对任意正整数 \(p\)，有

\[
      \sum_{k=1}^{p}a_kb_k
      =
      a_pB_p-\sum_{k=1}^{p-1}(a_{k+1}-a_k)B_k.
    \]

这也叫分部求和公式。它是离散版的分部积分：\(\sum a_kb_k\) 类比 \(\int f\,dg\)，部分和 \(B_k\) 类比原函数。

\subsubsection*{引理 9.4.2：Abel 引理}

若 \(\{a_k\}\) 单调，且 \(\{B_k\}\) 有界，即存在 \(M>0\)，使得对一切 \(k\)，

\[
      |B_k|\le M,
    \]

则

\[
      \left|\sum_{k=1}^{p}a_kb_k\right|
      \le M\left(|a_1|+2|a_p|\right).
    \]

直观上：如果 \(b_k\) 的部分和不乱跑，而 \(a_k\) 单调，那么乘起来后的部分和也能被控制。

\subsubsection*{定理 9.4.3：级数的 A-D 判别法}

若下列两个条件之一满足，则级数 \(\sum_{n=1}^{\infty}a_nb_n\) 收敛：

\begin{enumerate}
 \item \textbf{Abel 判别法：}\(\{a_n\}\) 单调有界，且 \(\sum_{n=1}^{\infty}b_n\) 收敛；
 \item \textbf{Dirichlet 判别法：}\(\{a_n\}\) 单调趋于 \(0\)，且 \(\left\{\sum_{i=1}^{n}b_i\right\}\) 有界。
 \end{enumerate}

\begin{center}
\small
\begin{tabularx}{\linewidth}{|>{\raggedright\arraybackslash}p{0.18\linewidth}|>{\raggedright\arraybackslash}p{0.42\linewidth}|>{\raggedright\arraybackslash}X|}
\hline
\textbf{类型} & \textbf{怎么认} & \textbf{常见题型} \\ \hline
Leibniz & \((-1)^{n+1}u_n\)，且 \(u_n\downarrow0\) & 交错调和、交错 \(p\)-型 \\ \hline
Dirichlet & 一个因子单调趋 \(0\)，另一个因子的部分和有界 & \(\sum a_n\sin nx\)、\(\sum a_n\cos nx\) \\ \hline
Abel & 一个因子单调有界，另一个因子本身成收敛级数 & 已知 \(\sum b_n\) 收敛，讨论 \(\sum a_nb_n\) \\ \hline
\end{tabularx}
\end{center}

\subsubsection*{例：三角级数}

若 \(\{a_n\}\) 单调趋于 \(0\)，则对任意实数 \(x\)，级数

\[
      \sum_{n=1}^{\infty}a_n\sin nx
    \]

收敛。理由是：当 \(x\ne 2k\pi\) 时，\(\sum_{k=1}^{n}\sin kx\) 有界；当 \(x=2k\pi\) 时，所有项为 \(0\)。因此由 Dirichlet 判别法收敛。

\subsection*{四、绝对收敛与条件收敛}

\subsubsection*{定义 9.4.2}

若

\[
      \sum_{n=1}^{\infty}|x_n|
    \]

收敛，则称 \(\sum x_n\) 为绝对收敛级数。若 \(\sum x_n\) 收敛而 \(\sum |x_n|\) 发散，则称 \(\sum x_n\) 为条件收敛级数。

\subsubsection*{基本事实}

\[
      \sum |x_n|\text{ 收敛}\quad\Longrightarrow\quad \sum x_n\text{ 收敛}.
    \]

证明入口是 Cauchy 收敛原理与三角不等式：

\[
      \left|x_{n+1}+\cdots+x_m\right|
      \le
      |x_{n+1}|+\cdots+|x_m|.
    \]

\subsubsection*{定理 9.4.4：正项和负项的拆分}

对任意项 \(x_n\)，定义

\[
      x_n^+=\frac{|x_n|+x_n}{2},\qquad
      x_n^-=\frac{|x_n|-x_n}{2}.
    \]

则

\[
      x_n=x_n^+-x_n^-,\qquad |x_n|=x_n^++x_n^-.
    \]

若 \(\sum x_n\) 绝对收敛，则 \(\sum x_n^+\) 与 \(\sum x_n^-\) 都收敛；若 \(\sum x_n\) 条件收敛，则 \(\sum x_n^+\) 与 \(\sum x_n^-\) 都发散到 \(+\infty\)。

\subsubsection*{例：幂参数级数}

讨论

\[
      \sum_{n=1}^{\infty}\frac{x^n}{n^p}
    \]

的敛散性：

\begin{center}
\small
\begin{tabularx}{\linewidth}{|>{\raggedright\arraybackslash}p{0.18\linewidth}|>{\raggedright\arraybackslash}p{0.42\linewidth}|>{\raggedright\arraybackslash}X|}
\hline
\textbf{参数} & \textbf{结论} & \textbf{依据} \\ \hline
\(|x|<1\) & 绝对收敛 & 根值判别，根值极限为 \(|x|\) \\ \hline
\(|x|>1\) & 发散 & 一般项不趋于 \(0\) \\ \hline
\(x=1,\ p>1\) & 绝对收敛 & \(p\)-级数 \\ \hline
\(x=1,\ p\le1\) & 发散 & \(p\)-级数 \\ \hline
\(x=-1,\ p>1\) & 绝对收敛 & \(\sum 1/n^p\) 收敛 \\ \hline
\(x=-1,\ 0<p\le1\) & 条件收敛 & Leibniz 收敛，但绝对值级数发散 \\ \hline
\(x=-1,\ p\le0\) & 发散 & 一般项不趋于 \(0\) \\ \hline
\end{tabularx}
\end{center}

\subsection*{五、加法交换律、重排与 Riemann 定理}

有限和可以随便换顺序，但无穷级数不一定。对收敛级数 \(\sum x_n\)，把所有项重新排列得到的新级数称为它的重排级数。

\subsubsection*{定理 9.4.5：绝对收敛级数可重排}

若 \(\sum x_n\) 绝对收敛，则任意重排级数 \(\sum x_n'\) 也绝对收敛，且和不变：

\[
      \sum_{n=1}^{\infty}x_n'
      =
      \sum_{n=1}^{\infty}x_n.
    \]

\subsubsection*{定理 9.4.6：Riemann 重排定理}

若 \(\sum x_n\) 条件收敛，则对任意给定的 \(a\)，其中 \(-\infty\le a\le+\infty\)，都存在 \(\sum x_n\) 的一个重排级数 \(\sum x_n'\)，使得

\[
      \sum_{n=1}^{\infty}x_n'=a.
    \]

直观做法是：正项加到刚超过 \(a\)，再加负项到刚低于 \(a\)，然后重复。因为正项和负项都分别发散到 \(+\infty\)，但每一项趋于 \(0\)，这种上下摆动会逐渐贴近 \(a\)。

\begin{notebox}
这说明“条件收敛”非常脆弱：它的收敛依赖原来的排列顺序。绝对收敛才真正保留了有限求和的交换律。
\end{notebox}

\subsection*{六、级数的乘法}

有限和相乘时，所有 \(a_ib_j\) 项加起来，顺序不影响结果。无穷级数相乘时，必须规定这些 \(a_ib_j\) 如何排列。

\subsubsection*{Cauchy 乘积}

给定两个级数 \(\sum a_n\)、\(\sum b_n\)，令

\[
      c_n=\sum_{i+j=n+1}a_ib_j
      =
      a_1b_n+a_2b_{n-1}+\cdots+a_nb_1.
    \]

则

\[
      \sum_{n=1}^{\infty}c_n
    \]

称为 \(\sum a_n\) 与 \(\sum b_n\) 的 Cauchy 乘积。

\subsubsection*{正方形排列与绝对收敛}

若按正方形方式排列所有 \(a_ib_j\)，则只要 \(\sum a_n\)、\(\sum b_n\)、以及正方形排列所得级数都收敛，就有

\[
      \sum d_n
      =
      \left(\sum_{n=1}^{\infty}a_n\right)
      \left(\sum_{n=1}^{\infty}b_n\right).
    \]

\subsubsection*{定理 9.4.7}

若 \(\sum a_n\) 与 \(\sum b_n\) 都绝对收敛，则把所有 \(a_ib_j\) 以任意方式排列求和所得的级数也绝对收敛，且其和等于

\[
      \left(\sum_{n=1}^{\infty}a_n\right)
      \left(\sum_{n=1}^{\infty}b_n\right).
    \]

\subsubsection*{典型应用：指数函数加法公式}

用 d'Alembert 判别法可知，对任意 \(x\in\mathbb R\)，

\[
      f(x)=\sum_{n=0}^{\infty}\frac{x^n}{n!}
    \]

绝对收敛。两个绝对收敛级数的 Cauchy 乘积给出：

\[
      \left(\sum_{n=0}^{\infty}\frac{x^n}{n!}\right)
      \left(\sum_{n=0}^{\infty}\frac{y^n}{n!}\right)
      =
      \sum_{n=0}^{\infty}\frac{(x+y)^n}{n!}.
    \]

也就是

\[
      f(x+y)=f(x)f(y).
    \]

\subsection*{七、判题路线图}

\begin{center}
\small
\begin{tabularx}{\linewidth}{|>{\raggedright\arraybackslash}p{0.18\linewidth}|>{\raggedright\arraybackslash}p{0.42\linewidth}|>{\raggedright\arraybackslash}X|}
\hline
\textbf{看到什么} & \textbf{先试什么} & \textbf{提醒} \\ \hline
\(\sum |x_n|\) 容易判 & 绝对收敛 & 一旦绝对收敛，原级数立刻收敛。 \\ \hline
\((-1)^n u_n\) & Leibniz & 检查 \(u_n\downarrow0\)，并说明是否绝对收敛。 \\ \hline
\(\sin nx,\cos nx\) & Dirichlet & 证明三角部分和有界，另一个因子单调趋零。 \\ \hline
已知 \(\sum b_n\) 收敛，再乘 \(a_n\) & Abel & 常见条件是 \(a_n\) 单调有界。 \\ \hline
重排或交换顺序 & 先问绝对收敛吗 & 条件收敛不能随便重排。 \\ \hline
乘法或 Cauchy 乘积 & 先看绝对收敛 & 两个都绝对收敛时最稳。 \\ \hline
\end{tabularx}
\end{center}

\subsection*{八、即时判断练习}

\begin{quickcheck}
\[
          \sum_{n=1}^{\infty}\frac{(-1)^{n+1}}{n}
        \]

\tcblower
\textbf{答案：}conditional
\end{quickcheck}

\begin{quickcheck}
\[
          \sum_{n=1}^{\infty}\frac{\sin nx}{n}\quad(0<x<2\pi)
        \]

\tcblower
\textbf{答案：}dirichlet
\end{quickcheck}

\begin{quickcheck}
\[
          \sum_{n=1}^{\infty}\frac{(-1)^n}{n^2}
        \]

\tcblower
\textbf{答案：}absolute
\end{quickcheck}

\subsection*{九、课后题}

下面按教材第九章第 4 节习题整理。题面完整保留，提示只给入口；写作业时仍建议先独立算一遍，再展开提示核对方向。

\begin{exercisebox}
\textbf{习题 1 \badge{敛散性、绝对收敛与条件收敛}}

讨论下列级数的敛散性，包括条件收敛与绝对收敛：

\[
      \begin{aligned}
      &(1)\ 1-\frac1{2!}+\frac13-\frac1{4!}+\frac15+\cdots;\\
      &(2)\ \sum_{n=1}^{\infty}\frac{(-1)^{n+1}}{n+x}\quad(x\ne -n);\\
      &(3)\ \sum_{n=1}^{\infty}(-1)^{n+1}\sin\frac{x}{n};\qquad
      (4)\ \sum_{n=1}^{\infty}\frac{(-1)^{n+1}}{\sqrt[n]{n}};\\
      &(5)\ \sum_{n=2}^{\infty}(-1)^n\frac{\ln^2 n}{n};\qquad
      (6)\ \sum_{n=1}^{\infty}\frac1{\sqrt n}\cos\frac{n\pi}{3};\\
      &(7)\ \sum_{n=1}^{\infty}(-1)^{n-1}\frac{4^n\sin^{2n}x}{n};\\
      &(8)\ \sum_{n=1}^{\infty}\frac{\sin(n+1)x\cos(n-1)x}{n^p};\\
      &(9)\ \sum_{n=1}^{\infty}(-1)^{n+1}\frac{n^2}{2^{x^n}};\\
      &(10)\ \sum_{n=1}^{\infty}(-1)^{n+1}
      \frac{\ln\left(2+\frac1n\right)}{\sqrt{(3n-2)(3n+2)}};\\
      &(11)\ \sum_{n=2}^{\infty}\frac{x^n}{n^p\ln^q n};\qquad
      (12)\ \sum_{n=1}^{\infty}\frac{(-1)^{n+1}}{n}\frac{a}{1+a^n}\quad(a>0).
      \end{aligned}
      \]

\begin{hintbox}
优先分三步：先看绝对值级数；再看是否 Leibniz；含 \(\sin nx,\cos nx\) 的题试 Dirichlet。第 (11) 是参数题，端点 \(x=\pm1\) 要单独讨论。
\end{hintbox}
\end{exercisebox}

\begin{exercisebox}
\textbf{习题 2 \badge{Cauchy 收敛原理判发散}}

利用 Cauchy 收敛原理证明下列级数发散：

\[
      \begin{aligned}
      &(1)\ 1+\frac12-\frac13+\frac14+\frac15-\frac16+\frac17+\frac18-\frac19+\cdots;\\
      &(2)\ 1-\frac12+\frac13+\frac14-\frac15+\frac16+\frac17-\frac18+\frac19+\cdots.
      \end{aligned}
      \]

\begin{hintbox}
找一段尾和，它不会随 \(n\) 增大而任意小。例如每三个一组或观察正项块与负项块的净贡献。
\end{hintbox}
\end{exercisebox}

\begin{exercisebox}
\textbf{习题 3 \badge{正项级数的尾部估计}}

设正项级数 \(\sum_{n=1}^{\infty}x_n\) 收敛，\(\{x_n\}\) 单调减少，利用 Cauchy 收敛原理证明：

\[
      \lim_{n\to\infty}nx_n=0.
      \]

\begin{hintbox}
对 \(x_n+x_{n+1}+\cdots+x_{2n}\) 用 Cauchy 收敛原理；再用单调性把它和 \(nx_{2n}\) 或 \(nx_n\) 联系起来。
\end{hintbox}
\end{exercisebox}

\begin{exercisebox}
\textbf{习题 4 \badge{Cauchy 条件的量词}}

若对任意 \(\varepsilon>0\) 和任意正整数 \(p\)，存在 \(N(\varepsilon,p)\)，使得

\[
      \left|x_{n+1}+x_{n+2}+\cdots+x_{n+p}\right|<\varepsilon
      \]

对一切 \(n>N\) 成立，问级数 \(\sum_{n=1}^{\infty}x_n\) 是否收敛？

\begin{hintbox}
注意 \(N\) 依赖 \(p\)。Cauchy 收敛原理要求一个 \(N(\varepsilon)\) 同时控制所有 \(p\)。试用调和级数作反例。
\end{hintbox}
\end{exercisebox}

\begin{exercisebox}
\textbf{习题 5 \badge{等价一般项}}

若级数 \(\sum_{n=1}^{\infty}x_n\) 收敛，且

\[
      \lim_{n\to\infty}\frac{x_n}{y_n}=1,
      \]

问级数 \(\sum_{n=1}^{\infty}y_n\) 是否收敛？

\begin{hintbox}
在正项级数中可以用极限比较。若是任意项级数，需要额外小心：等价一般项不能随便保留条件收敛。
\end{hintbox}
\end{exercisebox}

\begin{exercisebox}
\textbf{习题 6 \badge{交错级数的单调条件}}

设 \(x_n\ge0\)，且 \(\lim_{n\to\infty}x_n=0\)，问交错级数

\[
      \sum_{n=1}^{\infty}(-1)^{n+1}x_n
      \]

是否收敛？

\begin{hintbox}
Leibniz 判别法还需要单调减少。构造反例时，可以让奇数项像 \(1/k\)，偶数项很小。
\end{hintbox}
\end{exercisebox}

\begin{exercisebox}
\textbf{习题 7 \badge{由交错发散推出另一级数收敛}}

设正项数列 \(\{x_n\}\) 单调减少，且级数 \(\sum_{n=1}^{\infty}(-1)^nx_n\) 发散。问级数

\[
      \sum_{n=1}^{\infty}\left(\frac1{1+x_n}\right)^n
      \]

是否收敛？并说明理由。

\begin{hintbox}
若交错级数在单调正项条件下仍发散，说明 \(x_n\) 不趋于 \(0\)。于是 \(x_n\) 有正下界，后面的级数可与几何级数比较。
\end{hintbox}
\end{exercisebox}

\begin{exercisebox}
\textbf{习题 8 \badge{Abel 判别法}}

设级数 \(\sum_{n=1}^{\infty}\frac{x_n}{n^{\alpha_0}}\) 收敛，则当 \(\alpha>\alpha_0\) 时，级数

\[
      \sum_{n=1}^{\infty}\frac{x_n}{n^\alpha}
      \]

也收敛。

\begin{hintbox}
写成 \(\frac{x_n}{n^\alpha}=\frac1{n^{\alpha-\alpha_0}}\cdot\frac{x_n}{n^{\alpha_0}}\)，再用 Abel 判别法。
\end{hintbox}
\end{exercisebox}

\begin{exercisebox}
\textbf{习题 9 \badge{分部求和与差分}}

若 \(\sum_{n=1}^{\infty}nx_n\) 收敛，且 \(\sum_{n=2}^{\infty}n(x_n-x_{n-1})\) 收敛，则级数

\[
      \sum_{n=1}^{\infty}x_n
      \]

收敛。

\begin{hintbox}
令 \(b_n=nx_n\)。把 \(n(x_n-x_{n-1})\) 改写成 \(b_n-b_{n-1}-x_{n-1}\)，再看部分和。
\end{hintbox}
\end{exercisebox}

\begin{exercisebox}
\textbf{习题 10 \badge{有界变差乘收敛级数}}

若 \(\sum_{n=2}^{\infty}(x_n-x_{n-1})\) 绝对收敛，\(\sum_{n=1}^{\infty}y_n\) 收敛，则级数

\[
      \sum_{n=1}^{\infty}x_ny_n
      \]

收敛。

\begin{hintbox}
\(\sum |x_n-x_{n-1}|\) 收敛说明 \(\{x_n\}\) 有界变差，从而有极限且有界。用 Abel 变换控制尾和。
\end{hintbox}
\end{exercisebox}

\begin{exercisebox}
\textbf{习题 11 \badge{Taylor 估计}}

设 \(f(x)\) 在 \([-1,1]\) 上具有二阶连续导数，且

\[
      \lim_{x\to0}\frac{f(x)}{x}=0.
      \]

证明级数

\[
      \sum_{n=1}^{\infty}f\left(\frac1n\right)
      \]

绝对收敛。

\begin{hintbox}
由条件推出 \(f(0)=0\)、\(f'(0)=0\)。再用 Taylor 公式得到 \(f(x)=O(x^2)\)。
\end{hintbox}
\end{exercisebox}

\begin{exercisebox}
\textbf{习题 12 \badge{Abel 判别法的逆向使用}}

已知任意项级数 \(\sum_{n=1}^{\infty}x_n\) 发散，证明级数

\[
      \sum_{n=1}^{\infty}\left(1+\frac1n\right)^nx_n
      \]

也发散。

\begin{hintbox}
反证。若变换后的级数收敛，再乘上 \(\left(1+\frac1n\right)^{-n}\)，这个因子单调有界并趋于 \(e^{-1}\)，可由 Abel 判别法推出 \(\sum x_n\) 收敛，矛盾。
\end{hintbox}
\end{exercisebox}

\begin{exercisebox}
\textbf{习题 13 \badge{比值条件推出 Leibniz 条件}}

设 \(x_n>0\)，且

\[
      \lim_{n\to\infty}n\left(\frac{x_n}{x_{n+1}}-1\right)>0.
      \]

证明：交错级数

\[
      \sum_{n=1}^{\infty}(-1)^{n+1}x_n
      \]

收敛。

\begin{hintbox}
条件说明从某项起 \(\frac{x_n}{x_{n+1}}>1+\frac{c}{n}\)，于是 \(x_n\) 最终单调减少，且可以证明 \(x_n\to0\)。再用 Leibniz 判别法。
\end{hintbox}
\end{exercisebox}

\begin{exercisebox}
\textbf{习题 14 \badge{重排后的交错调和级数}}

利用

\[
      1+\frac12+\frac13+\cdots+\frac1n-\ln n\to \gamma\qquad(n\to\infty),
      \]

其中 \(\gamma\) 是 Euler 常数，求下述 \(\sum_{n=1}^{\infty}\frac{(-1)^{n+1}}{n}\) 的重排级数的和：

\[
      1+\frac13-\frac12+\frac15+\frac17-\frac14+\frac19+\frac1{11}-\frac16+\cdots.
      \]

\begin{hintbox}
每两个正的奇数项后放一个负的偶数项。设第 \(n\) 组部分和为 \(\sum_{k=1}^n\left(\frac1{4k-3}+\frac1{4k-1}-\frac1{2k}\right)\)，用 \(H_n-\ln n\to\gamma\) 化简。
\end{hintbox}
\end{exercisebox}

\begin{exercisebox}
\textbf{习题 15 \badge{Cauchy 乘积}}

利用级数的 Cauchy 乘积证明：

\[
      (1)\quad
      \sum_{n=0}^{\infty}\frac1{n!}\cdot
      \sum_{n=0}^{\infty}\frac{(-1)^n}{n!}=1;
      \]

\[
      (2)\quad
      \left(\sum_{n=0}^{\infty}q^n\right)
      \left(\sum_{n=0}^{\infty}q^n\right)
      =
      \sum_{n=0}^{\infty}(n+1)q^n
      =
      \frac1{(1-q)^2}\qquad(|q|<1).
      \]

\begin{hintbox}
第 (1) 可看作 \(e\cdot e^{-1}=1\)，也可直接计算 Cauchy 乘积中 \(\sum_{k=0}^n\frac{(-1)^k}{k!(n-k)!}\)。第 (2) 的 Cauchy 乘积每个 \(q^n\) 出现 \(n+1\) 次。
\end{hintbox}
\end{exercisebox}

来源参考：陈纪修、于崇华、金路《数学分析》第三版下册，第九章 §4。建议学习时把本节和 §3 正项级数 连起来看：绝对收敛就是把任意项级数重新交给正项级数判别法处理；条件收敛才需要 Leibniz、Dirichlet、Abel 和重排思想。

如果你做题后要批改，先发已做部分即可；我会按“结论、判别法、关键条件是否满足”逐题检查。
